\paragraph{\texorpdfstring{$(y,\delta)$}{(y,delta)} 平面闭包模型：五节点、伴随二次与谱坐标反演}
沿用结论 \ref{con:xi-terminal-zm-leyang-abel-jacobi-ode} 的椭圆曲线
\[
E:\quad Y^{2}=X^{3}-X+\frac14,
\qquad
y:=X^{2}+Y-\frac12\in\QQ(E),
\qquad
\omega:=\frac{\dd X}{2Y}.
\]
令 Abel--Jacobi 坐标 \(u\) 满足 \(\dd u=\omega\)，并令
\[
\delta:=\frac{\dd y}{\dd u}\in\QQ(E).
\]
由 \(\dd y=(4XY+3X^{2}-1)\,\omega\)（结论 \ref{con:xi-terminal-zm-leyang-abel-jacobi-ode} 的证明中已给出）可得
\[
\delta=4XY+3X^{2}-1.
\]

\begin{theorem}[\texorpdfstring{$\delta$}{delta}-闭包的平面五次与五节点归一化]\label{thm:xi-terminal-zm-delta-closure-plane-quintic}
令
\[
F_\delta(T;y):=T^{4}+(8y-3)T^{3}+(2y^{2}-34y-4)T^{2}-y(y-1)\,P_{\mathrm{LY}}(y)\in\ZZ[y][T],
\]
\[
P_{\mathrm{LY}}(y):=256y^{3}+411y^{2}+165y+32.
\]
记仿射平面曲线
\[
\mathcal C_\delta:=\{(y,T)\in\mathbb A^{2}_{(y,T)}:\ F_\delta(T;y)=0\},
\]
并记 \(\overline{\mathcal C}_\delta\subset\mathbb P^{2}\) 为将 \(F_\delta\) 以总次数 \(5\) 齐次化所得的射影闭包。
则成立：
\begin{enumerate}
  \item \(\overline{\mathcal C}_\delta\) 在仿射部分 \(\mathbb A^{2}\) 内恰有五个奇点，并且全为普通双点（节点）；
  \item \(\overline{\mathcal C}_\delta\) 在无穷远处仅有一个点且为光滑点；
  \item \(\overline{\mathcal C}_\delta\) 的归一化为亏格 \(1\) 的光滑射影曲线，并与 \(E\) 在 \(\QQ\) 上同构；
  等价地，
  \[
  \QQ(\overline{\mathcal C}_\delta)\cong \QQ(y,T)/(F_\delta)\cong \QQ(E)=\QQ(X,Y).
  \]
\end{enumerate}
\end{theorem}

\begin{proof}[证明要点]
首先，由结论 \ref{con:xi-terminal-zm-leyang-abel-jacobi-ode} 得 \(\QQ(E)=\QQ(y,\delta)\)，并且 \((y,\delta)\) 满足 \(F_\delta(\delta;y)=0\)；
因此 \(\QQ(y,T)/(F_\delta)\) 的光滑模型为亏格 \(1\) 的曲线，并与 \(E\) 同构（函数域同构确定归一化）。

其次，齐次化取总次数 \(5\)，可写为
\[
F_\delta^{\mathrm{h}}(T,y,Z)=
Z\Bigl(T^{4}+(8y-3Z)T^{3}+(2y^{2}-34yZ-4Z^{2})T^{2}\Bigr)-y(y-Z)P_{\mathrm{LY}}^{\mathrm{h}}(y,Z),
\]
其中
\(
P_{\mathrm{LY}}^{\mathrm{h}}(y,Z)=256y^{3}+411y^{2}Z+165yZ^{2}+32Z^{3}
\)。
令 \(Z=0\) 得 \(F_\delta^{\mathrm{h}}(T,y,0)=-256y^{5}\)，故无穷远处唯一点为 \([T:y:Z]=[1:0:0]\)。
并且
\(
\partial_ZF_\delta^{\mathrm{h}}(1,0,0)=1\neq 0
\)，因此该无穷远点光滑，得到第二条。

对仿射奇点，设
\[
I_{\mathrm{sing}}:=(F_\delta,\partial_TF_\delta,\partial_yF_\delta)\subset\QQ[T,y].
\]
由于奇点必满足 \(F_\delta(T;y)\) 在 \(T\) 上有重根，故其 \(y\)-坐标只能落在判别式零点上。
由结论 \ref{con:xi-terminal-zm-leyang-abel-jacobi-ode} 的判别式分解
\[
\mathrm{Disc}_{T}(F_\delta)=-y(y-1)P_{\mathrm{LY}}(y)\,Q_{5}(y)^{2},
\quad
Q_{5}(y)=4096y^{5}+5376y^{4}-464y^{3}-2749y^{2}-723y+80,
\]
奇点候选仅可能出现在 \(y\in\{0,1\}\cup\{P_{\mathrm{LY}}(y)=0\}\cup\{Q_5(y)=0\}\)。

当 \(T=0\) 时，\(F_\delta(0;y)=-y(y-1)P_{\mathrm{LY}}(y)\)。
直接计算
\[
\partial_yF_\delta(0;y)=-(2y-1)P_{\mathrm{LY}}(y)-y(y-1)P_{\mathrm{LY}}'(y),
\]
并注意 \(P_{\mathrm{LY}}\) 无重根且其根不等于 \(0,1\)，可知在 \(y=0,1\) 与 \(P_{\mathrm{LY}}(y)=0\) 的各点处 \(\partial_yF_\delta(0;y)\neq 0\)；
因此 \((y,T)=(0,0)\)、\((1,0)\) 以及 \(\{P_{\mathrm{LY}}=0\}\times\{0\}\) 仅对应分歧而非平面奇点。
从而一切平面奇点均满足 \(T\neq 0\)，并且其 \(y\)-坐标必满足 \(Q_5(y)=0\)。

在 \(I_{\mathrm{sing}}\) 上可作消元得到
\(
I_{\mathrm{sing}}\cap\QQ[y]=(Q_5(y))
\)，并且在 \(\QQ[T,y]/I_{\mathrm{sing}}\) 中有线性关系
\[
93\,T=P_4(y),
\qquad
P_4(y)=8192y^{4}+5888y^{3}-3680y^{2}-2848y+152.
\]
由于 \(\gcd(Q_5,Q_5')=1\)，故 \(Q_5\) 的五个根互异，从而得到五个互异奇点
\[
\left(y,\ T=\frac{P_4(y)}{93}\right),\qquad Q_5(y)=0.
\]

最后，为判定奇点类型，取二次切锥判别式
\[
\Delta:=(\partial_{Ty}F_\delta)^{2}-(\partial_{TT}F_\delta)(\partial_{yy}F_\delta).
\]
将 \(\Delta\) 在 \(\QQ[T,y]/I_{\mathrm{sing}}\) 中约化得到
\[
\Delta\equiv 2\cdot\bigl(19072y^{4}+51120y^{3}+45129y^{2}+14759y+1032\bigr).
\]
该四次多项式与 \(Q_5\) 互素，故在 \(Q_5(y)=0\) 的各点处 \(\Delta\neq 0\)，从而切锥分解为两条不同直线，五个奇点均为节点。
由于平面五次曲线的算术亏格为 \(\frac{(5-1)(5-2)}{2}=6\)，节点的 \(\delta\)-不变量均为 \(1\)，总和为 \(5\)，故几何亏格为 \(1\)。
与第一段的函数域同构合并，即得第三条。证毕。
\end{proof}

\begin{proposition}[伴随二次曲线的唯一性与 \texorpdfstring{$Q_5$}{Q5} 的内禀消元刻画]\label{prop:xi-terminal-zm-delta-adjoint-conic-and-elimination}
令
\[
H(T,y):=\frac{1}{T}\,\partial_TF_\delta(T;y)=4T^{2}+(24y-9)T+(4y^{2}-68y-8).
\]
则：
\begin{enumerate}
  \item \(H(T,y)=0\) 在射影意义下为 \(\overline{\mathcal C}_\delta\) 的唯一伴随二次曲线（即通过全部五个节点的二次曲线，按比例唯一）；
  \item 结果式消元满足恒等式
  \[
  \mathrm{Res}_{T}\bigl(F_\delta(T;y),\,H(T,y)\bigr)=Q_{5}(y)^{2}.
  \]
  因而 \(Q_5\) 可被内禀刻画为理想 \((F_\delta,\frac{1}{T}\partial_TF_\delta)\) 的消元因子。
\end{enumerate}
\end{proposition}

\begin{proof}[证明要点]
由于 \(F_\delta\) 不含 \(T\) 的一次项，\(\partial_TF_\delta\) 含因子 \(T\)，故 \(H\) 为二次多项式。
任一节点满足 \(F_\delta=\partial_TF_\delta=\partial_yF_\delta=0\)，并且由定理 \ref{thm:xi-terminal-zm-delta-closure-plane-quintic} 的证明知节点处 \(T\neq 0\)，故节点亦满足 \(H=0\)；
于是 \(H=0\) 通过全部五个节点。

另一方面，次数为 \(5\) 且仅含节点的平面曲线，其正则微分由次数 \(d-3=2\) 的伴随曲线给出；
在亏格 \(1\) 情形，伴随二次曲线空间维数为 \(1\)，故按比例唯一。
因此 \(H=0\) 为唯一伴随二次曲线。

对结果式，使用判别式--结果式恒等式
\[
\mathrm{Disc}_{T}(F_\delta)=\pm\,\mathrm{Res}_{T}(F_\delta,\partial_TF_\delta),
\]
并用 \(\partial_TF_\delta=T\cdot H\) 与结果式乘法性得到
\[
\mathrm{Res}_{T}(F_\delta,\partial_TF_\delta)
=\mathrm{Res}_{T}(F_\delta,T)\,\mathrm{Res}_{T}(F_\delta,H).
\]
注意 \(\mathrm{Res}_{T}(F_\delta,T)=F_\delta(0;y)=-y(y-1)P_{\mathrm{LY}}(y)\)。
与结论 \ref{con:xi-terminal-zm-leyang-abel-jacobi-ode} 的判别式分解合并，即得
\(
\mathrm{Res}_{T}(F_\delta,H)=Q_5(y)^2
\)。
证毕。
\end{proof}

\begin{theorem}[谱坐标 \texorpdfstring{$\lambda=X$}{lambda=X} 的有理反演与极点除子]\label{thm:xi-terminal-zm-delta-spectral-x-inversion}
继续沿用上述记号，并令谱坐标 \(\lambda:=X\)。
在开集
\[
U:=\{(y,T)\in \mathcal C_\delta:\ Q_{5}(y)\neq 0\}
\]
上，有显式有理反演
\[
\lambda=\frac{N(y,T)}{Q_{5}(y)},
\]
其中
\[
\begin{aligned}
N(y,T)=\;&32T^{3}y+19T^{3}+256T^{2}y^{3}+416T^{2}y^{2}+24T^{2}y-81T^{2}\\
&+512Ty^{4}+160Ty^{3}-709Ty^{2}-411Ty-20T\\
&-1792y^{5}-2528y^{4}-914y^{3}-377y^{2}-85y+80.
\end{aligned}
\]
并令
\[
Y:=y-\lambda^{2}+\frac12.
\]
则 \((\lambda,Y)\) 满足椭圆曲线方程
\[
Y^{2}=\lambda^{3}-\lambda+\frac14,
\]
且反演关系
\[
y=\lambda^{2}+Y-\frac12,
\qquad
T=4\lambda Y+3\lambda^{2}-1
\]
成立。
因此在 \(U\) 上，映射
\[
(\lambda,Y)\longmapsto (y,T)=\Bigl(\lambda^{2}+Y-\frac12,\ 4\lambda Y+3\lambda^{2}-1\Bigr)
\]
给出 \(E\dashrightarrow\mathcal C_\delta\) 的双有理参数化，而 \(\lambda=N/Q_{5}\) 给出其在 \(U\) 上的显式逆。
\end{theorem}

\begin{proof}[证明要点]
由结论 \ref{con:xi-terminal-zm-leyang-abel-jacobi-ode} 的消元链条，\(\QQ(E)=\QQ(y,T)\)，故 \(\lambda\in\QQ(y,T)\)。
另一方面，\((\lambda,y)\) 必满足谱四次
\[
\Pi(\lambda,y)=\lambda^{4}-\lambda^{3}-(2y+1)\lambda^{2}+\lambda+y(y+1)=0,
\]
并且由 \(T=4\lambda Y+3\lambda^{2}-1\) 与 \(Y=y-\lambda^{2}+\frac12\) 消去 \(Y\) 得到三次约束
\[
4\lambda^{3}-3\lambda^{2}-(4y+2)\lambda+(1+T)=0.
\]
在 \(\QQ[\lambda,y,T]\) 中对上述两式作消元，可得到线性关系
\[
Q_{5}(y)\lambda-N(y,T)=0,
\]
其中 \(Q_5\) 与 \(N\) 如上。
当 \(Q_5(y)\neq 0\) 时可解得 \(\lambda=N/Q_5\)。

取 \(Y:=y-\lambda^2+\tfrac12\)，则由定义有 \(y=\lambda^2+Y-\tfrac12\)。
将 \(\Pi(\lambda,y)=0\) 代入并整理即得
\(
Y^2=\lambda^3-\lambda+\tfrac14
\)，从而 \((\lambda,Y)\in E\)。
最后将 \((\lambda,Y)\) 代回 \(T=4\lambda Y+3\lambda^2-1\) 得恒等成立，故得双有理互逆。证毕。
\end{proof}

\begin{corollary}[判别式分解中平方因子的几何来源]\label{cor:xi-terminal-zm-delta-discriminant-square-geometry}
令 \(\widetilde{\mathcal C}_\delta\to\overline{\mathcal C}_\delta\) 为归一化，\(\pi_y:\widetilde{\mathcal C}_\delta\to\mathbb P^1_y\) 为由坐标 \(y\) 诱导的投影。
则：
\begin{enumerate}
  \item 因子 \(-y(y-1)P_{\mathrm{LY}}(y)\) 恰给出 \(\pi_y\) 的分歧像（有限分歧值），对应于 \(T=0\) 的分支；
  \item 多项式 \(Q_{5}(y)\) 的零点不对应 \(\pi_y\) 的分歧值，而对应于平面闭包 \(\overline{\mathcal C}_\delta\subset\mathbb P^{2}\) 的五个节点在 \(y\)-直线上的投影；
  \item 因而
  \[
  \mathrm{Disc}_{T}(F_\delta(T;y))
  =-\underbrace{y(y-1)P_{\mathrm{LY}}(y)}_{\text{归一化覆盖的分歧贡献}}
  \cdot
  \underbrace{Q_{5}(y)^{2}}_{\text{平面模型节点造成的平方因子}}
  \]
  给出“分歧贡献与平面自交贡献”的正交分解：平方因子完全记录平面化表示中的节点识别，而不改变归一化后的本征分歧集合。
\end{enumerate}
\end{corollary}

\begin{proof}[证明要点]
由命题 \ref{prop:xi-terminal-zm-delta-adjoint-conic-and-elimination} 的结果式分解
\[
\mathrm{Res}_{T}(F_\delta,\partial_TF_\delta)
=\mathrm{Res}_{T}(F_\delta,T)\,\mathrm{Res}_{T}(F_\delta,H),
\]
其中 \(\mathrm{Res}_{T}(F_\delta,T)=F_\delta(0;y)=-y(y-1)P_{\mathrm{LY}}(y)\) 对应 \(T=0\) 处的重根，即 \(\pi_y\) 的分歧像；
而 \(\mathrm{Res}_{T}(F_\delta,H)=Q_5(y)^2\) 对应 \(T\neq 0\) 的约化导数零点。
由定理 \ref{thm:xi-terminal-zm-delta-closure-plane-quintic}，这些点恰为平面模型的节点，其 \(y\)-坐标由 \(Q_5(y)=0\) 给出。
因此平方因子仅记录平面闭包的自交，并不引入新的归一化分歧值。证毕。
\end{proof}

\endinput

